\section{From marked masks to the profile}\label{app:transfer}
Two squares can depend on the survival of the same line. Treating their indicators as independent
would change the pair mean, even if the profile of each square were exact. We therefore transfer the
weighted contribution of each state with its own mass. For a fixed pattern we identify equal affine
labels and take each distinct indicator only once. With \(d_*\) free indices, an accidental
coincidence satisfies one additional nonconstant affine equation and costs \(O(n^{-1})\) after
division by \(n^{d_*}\). On the auxiliary regular trajectory, at \(x\ge b\), the hazards for surviving
axial and diagonal lines are \(1/m\) and \(\ell_f/m\), respectively. The latter differ from their
deterministic values by at most \(\eta/n\) at fixed precision. Two distinct lines are occupied
simultaneously with probability \(O_b(n^{-2})\). The Bonferroni inequalities give the product of their
one-step survival probabilities with error \(O_{b,v}(\eta/n+n^{-2})\). We multiply each current
indicator by its deterministic future survival probability up to the cut-off. One piece changes
\(O_v(n^{d_*-1})\) summands. Azuma's inequality for the centered mean, with jumps bounded by
\(O_{b,v}(1/n)\), and the regularity bound give an exponential tail at fixed precision. The union over
marked lines has polynomial size. The density ratios up to the stopping time preserve the tail,
including the bound \(\exp(O_b(n^{2/3}))\) for \(X\) when changing to \(\mu\). Thus each surviving
axial label contributes \(b\) and each distinct diagonal label contributes \(U_f(b)\) to the
normalized mean.

In particular, summing \(YW_b^0/j\) along each surviving line gives margins \(1+o_{\Pr}(1)\)
uniformly, and two surviving columns have \(\ge c_bj\) common neighbors. In logarithmic column
potentials, the map that normalizes first rows and then columns has a stochastic derivative with
entries \(\ge c'_b/j\) on the unit ball in the oscillation seminorm. Its contraction in oscillation
gives \(jD_{il}=Y_{il}W_b^0(i,l)(1+o_{\Pr}(1))\) for the Sinkhorn scaling, uniformly on the support.
When the two distinguished squares are removed, each margin loses \(O_b(1/j)\) and two columns lose at
most two common neighbors. The same proof still applies. The stationary rook occupancy satisfies
\(p_{il}/D_{il}=1+O_b(1/j)\) on the support. In the integrated patterns, the incidences from an
additional row \(y\) to the marked columns have a product law with parameters
\(v_b(y,c_0),v_b(y,c_1)\). Distinct columns give distinct diagonal labels in each family, except for
the \(O(1)\) values of \(y\) that coincide with a distinguished label. We may multiply the mean by any
of the four bounded state indicators. For the quadratic term of \(\eta\), let \(B=p_b-1\), \(G_R=BB^*,
T_R=G_R(I-G_R)^{-1}, U=T_RB\), with the transposed operators for columns. Let
\(d_R=\operatorname{diag}(T_R-G_R)\) and let \(d_C\) be its column analogue, with \(\bar d=\int
d_R=\int d_C\). Define \(\psi_R=2[(B+U)d_C-T_Rd_R]-d_R\) and its column counterpart. At the four
marked endpoints (two rows \(r_0,r_1\) and two columns \(c_0,c_1\)), the derivative identities give
\begin{align*}
 R''_1(0)/2={}&XW^0_{01}+YW^0_{10}+\sum_{\square} U-T_R(r_0,r_1)-T_C(c_0,c_1)
              -\tfrac12\sum_* d,\\
 A'_1(0)={}&-1+2\bar d+\tfrac12\sum_* \psi .
\end{align*}
The asterisk ranges over the four endpoints and the square symbol over its four board squares, with
the columns ordered so that the distinguished squares are \(00,11\). For the first identity we
differentiate the margins with respect to the occupancy parameter. On the ordinary types the mean
response solves \((I+N)y=N\alpha\), where \(N\) has off-diagonal blocks \(p_b,p_b^*\) and \(\alpha\)
inserts the two endpoints on each side. The centered blocks of \(N(I+N)^{-1}N\) are
\(T_R,-U,-U^*,T_C\). At each marked vertex a repeated incidence contributes its Bernoulli variance.
For the prefactor we differentiate the local factor, the margin diagonals, the total mass and the
centered determinant. Their sum is \(1+\langle d-1,y\rangle-\sum_*d/2\), and substituting the centered
solution and its constant mode gives the second identity. The limiting linear part \(jD\) is the sum
\(L_s\) in each submeasure. Squaring gives \(L_s^2/2\).

The pattern argument applies at any fixed length, but the second coefficient contains the inverse of
the margin matrix and hence walks of arbitrary length. We truncate its expansion in centered
alternating walks. Each walk of length \(q\ge2\) between fixed endpoints is bounded by
\(C_b^2r_b^{q-2}\), by the two section norms and the gap \(r_b<1\). At second order in the logarithm
of the determinant there are at most two insertions in each trace, and the tail is bounded by
\(C_b\sum_{q>N}(1+q)^2r_b^{\max(2q-C,0)}\). At a fixed truncation, the margin recurrences attach
rooted trees to a square or an alternating trace. Coincidences of distinct free labels cost
\(O_{b,N}(1/n)\). Each repeated indicator is retained before averaging. A trace of length two on a
free edge with conditional mean \(V\) uses \(Y^2=Y\) on the outward and return steps. The excess over
the product of the means contributes \(V(1-V)\ell^2/(2d_i d_l)\), with \(\ell\) the normalized edge
weight in the Hessian and \(d_i,d_l\) its row and column degrees. The root edges already have their
indicators fixed. Thus the ordinary variance is counted once. First we take \(n\to\infty\) at fixed
truncations, then we remove the truncations. Only then do we take \(b\downarrow0\), since the
positivity of \(v_b\) gives a uniform lower bound for \(p_b^*p_b\) near zero and dominated convergence
in each integral. With \(\mathcal I_0(C)=\sum_o\int_\Omega w_{o,0}\sum_s q_{o,s,0}C_{o,s}\), we obtain
\[
 \lim_{b\downarrow0}\limsup_n\left|
 n\mathbb E_\mu[\textstyle\sum_{\{e,f\}}p_ep_f B_{\rm phys}(e,f);F_b]
 -\mathcal I_0(B_{2,s})\right|=0. \tag{F2}
\]
Thus the signed pair contribution at the boundary is \(J=-B_2+L_2\), with \(B_2=\mathcal
I_0(B_{2,s})\) and \(L_2=-\mathcal I_0(L_s^2/2)\). In state \(11\) each orientation contributes its
single-line coefficient, and the interaction of the two orientations appears only once in \(-S\).
Adding \(J\) to the other two means gives the constant \(P_T+J-S\) that enters the cut-off comparison
once the limit in \(n\) has been taken for each fixed cut-off.
